\documentclass[conference]{IEEEtran}
\IEEEoverridecommandlockouts

\usepackage{cite}
\usepackage{amsmath,amssymb,amsfonts}
\usepackage{algorithmic}
\usepackage{graphicx}
\usepackage{textcomp}
\usepackage{xcolor}
\usepackage{booktabs}
\usepackage{enumitem}
\usepackage{bm}        % \bm{} for bold math symbols
\usepackage{mathtools} % \coloneqq, dcases

% ─── Notation shorthands ─────────────────────────────────────────────────────
\newcommand{\R}{\mathbb{R}}
\newcommand{\Z}{\mathbb{Z}}
\newcommand{\OUV}{\Omega_{\text{UV}}}
\newcommand{\Oin}{\Omega_{\text{in}}}
\newcommand{\Osrc}{\Omega_{\text{src}}}
\newcommand{\Otgt}{\Omega_{\text{tgt}}}
\newcommand{\Vis}{\operatorname{Vis}}
\newcommand{\Tbaked}{T_{\text{baked}}}
\newcommand{\Tpartial}{T_{\text{partial}}}
\newcommand{\TIDM}{T_{\text{IDM}}}
\newcommand{\Tcomplete}{T_{\text{complete}}}
\newcommand{\Tout}{T_{\text{out}}}
\newcommand{\Vsrc}{\mathbf{V}_{\text{src}}}
\newcommand{\cM}{\mathcal{M}}       % validity mask
\newcommand{\cA}{\mathcal{A}}       % 4D avatar
% ─────────────────────────────────────────────────────────────────────────────

\def\BibTeX{{\rm B\kern-.05em{\sc i\kern-.025em b}\kern-.08em
    T\kern-.1667em\lower.7ex\hbox{E}\kern-.125emX}}

\begin{document}

\title{Complete-Texture 4D Face Avatar Reconstruction from a Single Image}

\author{
  \IEEEauthorblockN{[Author 1 Name]}
  \IEEEauthorblockA{\textit{[Department]} \\
    \textit{[University, Country]} \\
    author1@university.edu}
  \and
  \IEEEauthorblockN{[Author 2 Name]}
  \IEEEauthorblockA{\textit{[Department]} \\
    \textit{[University, Country]} \\
    author2@university.edu}
  \and
  \IEEEauthorblockN{[Author 3 Name]}
  \IEEEauthorblockA{\textit{[Department]} \\
    \textit{[University, Country]} \\
    author3@university.edu}
}

\maketitle

% ─────────────────────────────────────────────────────────────────────────────
\begin{abstract}
Monocular 3D face reconstruction methods recover accurate geometry from a
single image but yield only partial, camera-view-baked textures: occluded
surface regions---ears, scalp, and the back of the head---remain blank in the
UV atlas, making the reconstructed mesh unsuitable for free-viewpoint rendering
or 4D animated face synthesis. We present a complete pipeline that bridges this
gap by coupling a UV-space inpainting diffusion model with a geometry-aware
cross-topology texture transfer module, producing a full-surface texture atlas
that covers the entire head from a single photograph. The completed texture is
applied uniformly across an animated mesh sequence to form a 4D representation:
an avatar simultaneously navigable in three spatial dimensions and across the
temporal axis of expressive face motion. We formalize the 4D texture property
as a coverage condition on the UV domain and show that partial textures violate
this condition for the majority of the camera-pose sphere, while our complete
texture restores it unconditionally. Quantitative evaluation on a large
multi-view face dataset demonstrates consistent and substantial improvements
over the monocular reconstruction baseline across all evaluated subjects on
standard image quality metrics. Our method requires a single photograph as
input, generalizes across identities without per-subject training, and produces
a temporally coherent 4D avatar whose texture quality is uniform regardless of
viewing direction.
\end{abstract}

\begin{IEEEkeywords}
4D face avatar, complete texture, UV inpainting, cross-topology texture
transfer, free-viewpoint rendering, monocular face reconstruction
\end{IEEEkeywords}

% ─────────────────────────────────────────────────────────────────────────────
\section{Introduction}

The proliferation of remote communication---accelerated by global events that
forced billions of people to interact primarily through digital channels---has
exposed a fundamental gap between what video conferencing delivers and what
human presence truly requires. Flat 2D video streams are bandwidth-hungry
(1--5~Mbps for HD), privacy-invasive (the raw camera feed is transmitted
unconditionally), and viewpoint-locked: the receiver sees only the angle
captured by the sender's camera. More critically, 2D video cannot be
manipulated in 3D space, overlaid in augmented reality environments, or
integrated into spatial computing platforms such as virtual meeting rooms or
immersive e-learning systems.

A compelling alternative is \emph{parametric avatar streaming}: instead of
transmitting raw pixels, the sender's face is represented as a compact
parametric model whose expression parameters---roughly 500 floating-point
values per frame---are transmitted at a fraction of raw-video bandwidth.
Parameters compressed to 16-bit half-precision require under 1~KB per frame,
compared to approximately 30--150~KB per HD video frame. Since the receiver
renders the mesh locally, the avatar can be viewed from \emph{any angle},
enabling immersive free-viewpoint telepresence without additional hardware.

This vision hinges on a prerequisite that existing pipelines fail to satisfy:
\textbf{the avatar must possess a complete, photorealistic texture across its
entire surface.} Monocular 3D face reconstruction methods~\cite{feng2021deca}
recover accurate mesh geometry but bake the texture exclusively from the visible
pixels of the input photograph. Regions occluded at capture time---the back of
the head, ears, temples, and portions of the neck---remain blank in the UV
atlas. When the receiver rotates the avatar even slightly away from the original
capture angle, these gaps are immediately exposed, making free-viewpoint
rendering impractical (Fig.~1).

The need for complete-texture 3D avatars extends well beyond video
conferencing. In \textbf{e-learning and distance education}, a single
enrollment photograph can generate a persistent instructor avatar. In
\textbf{digital accessibility}, people with conditions affecting facial
appearance can adopt a fully-textured avatar. In \textbf{cultural and
historical preservation}, one archived photograph of a historical figure can
yield an interactive 3D memorial. In each scenario, texture completeness is the
decisive enabling factor.

UV-space inpainting models~\cite{bai2023uvidm} have demonstrated the ability
to synthesize plausible texture for occluded facial regions in a high-resolution
UV layout. However, a fundamental problem remains unsolved: the inpainting model
operates on its own mesh parameterization (e.g., BFM~\cite{paysan2009bfm}),
whereas the animation-ready reconstruction output uses a different mesh topology
(e.g., FLAME~\cite{li2017learning}). These two meshes differ in vertex count,
face connectivity, UV layout, and spatial coordinate convention. No prior work
provides a principled method to transfer a completed UV texture atlas across
this topology boundary while guaranteeing full surface coverage and no seam
artifacts---the gap this paper fills.

We address this with a cross-topology texture transfer algorithm that is
\emph{model-agnostic}: it treats source and target as arbitrary triangle meshes
and requires only vertex positions, face connectivity, UV coordinates, and
landmark correspondences. It therefore generalizes to any (source mesh,
target mesh) pair beyond the specific models used here.

\medskip\noindent\textbf{Contributions.} This paper makes the following
contributions:
\begin{itemize}[leftmargin=*,nosep]
  \item \textbf{Cross-topology UV texture transfer (primary contribution).}
    A general-purpose algorithm for transferring a UV texture atlas from a
    source mesh to an arbitrary target mesh with a different topology.
    The algorithm combines robust Procrustes~+~three-pass Point-to-Plane ICP
    alignment, soft-confidence vertex color transfer (guaranteeing 100\% surface
    coverage without hard distance gates), and a zero-loop vectorized UV
    rasterizer. Each design choice is independently validated by ablation.
    The algorithm is model-agnostic: it applies to any (source, target) mesh
    pair.
  \item \textbf{Formal 4D avatar definition.}
    A formal definition of the \emph{4D texture avatar} as a UV-domain coverage
    condition, establishing texture completeness as a necessary---not merely
    desirable---property for free-viewpoint animated face rendering, and
    quantifying the fraction of the camera-pose sphere that becomes inaccessible
    under a partial texture.
  \item \textbf{End-to-end pipeline and evaluation.}
    Integration of a monocular face reconstruction model and a UV-space
    inpainting model via the transfer module, producing a complete-texture 4D
    avatar from a single photograph in under 90~seconds. Quantitative evaluation
    on 74 subjects from the PolyFace dataset~\cite{polyface} shows consistent
    improvement over the partial-texture baseline across all subjects.
\end{itemize}

% ─────────────────────────────────────────────────────────────────────────────
\section{Related Work}

\medskip\noindent\textbf{Monocular 3D face reconstruction.}
3D Morphable Models~\cite{blanz1999morphable} provide a low-dimensional
parametric space for face shape and texture, enabling reconstruction from as
few as one image. Early fitting methods~\cite{paysan2009bfm} relied on
iterative optimization; recent deep learning approaches regress morphable model
parameters directly~\cite{tran2017regressing,deng2019accurate}.
State-of-the-art methods~\cite{feng2021deca,danecek2022emoca} extend this to
per-frame expression tracking, achieving high-quality reconstruction from
monocular video. All these methods share the same texture bottleneck: the
estimated albedo is reconstructed from the visible image region only, leaving a
large portion of the UV atlas blank or extrapolated poorly---the gap that
motivates this work.

\medskip\noindent\textbf{UV texture completion.}
Face texture completion in UV space has been addressed by GAN-based inpainting
models~\cite{deng2018uvgan} that learn to hallucinate realistic skin texture
for occluded regions conditioned on visible ones. More recently, diffusion model
backbones applied directly to UV atlas layouts~\cite{bai2023uvidm} achieve
significantly higher fidelity for complex skin tones and fine facial features.
All such models produce a texture atlas tied to the UV parameterization of their
training mesh (typically BFM~\cite{paysan2009bfm}). Animation-ready models
(e.g., FLAME~\cite{li2017learning}) use a different vertex count, face
connectivity, and UV layout. The completed texture therefore cannot be applied
to the animation mesh without a principled cross-topology transfer---which no
prior work provides. Naive solutions (nearest-vertex color copy, direct UV
remapping) fail because they ignore the geometric misalignment between the two
mesh coordinate systems and produce coverage gaps wherever the topology
structures diverge. Our algorithm is the first to address this transfer
problem explicitly.

\medskip\noindent\textbf{Animated face and 4D face synthesis.}
NeRF-based methods~\cite{guo2021adnerf,park2021nerfface} achieve
photo-realistic novel-view synthesis but are identity-specific, requiring
per-subject training and substantial compute at inference. Explicit mesh-based
animated face systems~\cite{thies2016face2face,fried2019text} operate in real
time and generalize across identities but suffer from the same
incomplete-texture limitation as other monocular methods. Speech-driven
animation models~\cite{xing2023codetalker} generate expressive face motion but
output 2D video, discarding the 3D representation entirely. Our approach is
complementary: we focus on completing the surface texture of a given animated
mesh sequence, restoring the free-viewpoint property that partial textures
violate.

% ─────────────────────────────────────────────────────────────────────────────
\section{Method}

\subsection{Pipeline Overview}

Our system transforms a monocular face video into a 4D avatar---a temporally
indexed sequence of complete-texture 3D meshes---through four successive stages
(Fig.~1):

\begin{enumerate}[leftmargin=*,nosep]
  \item \textbf{Face reconstruction and texture completion.} A monocular 3D
    face reconstruction method processes each video frame to obtain per-frame
    mesh geometry and a partial UV texture baked from the visible image region.
    A UV-space inpainting model then completes the full texture atlas from the
    identity frame, synthesizing plausible content for all occluded surface
    regions.
  \item \textbf{Cross-topology texture transfer.} A geometry-aware module that
    aligns the completed inpainting-mesh texture onto the target animation-mesh
    UV space via (i)~robust Procrustes alignment using 68 facial landmarks,
    (ii)~a three-pass ICP refinement scheme, (iii)~vectorized vertex-level color
    transfer with soft confidence weighting, (iv)~zero-loop UV rasterization and
    barycentric baking at $1024\!\times\!1024$ resolution, and (v)~seam dilation
    with Gaussian boundary blending.
  \item \textbf{4D assembly.} The completed texture---computed once per
    subject---is applied uniformly to every frame of the animation sequence,
    preserving per-frame expression dynamics in the mesh geometry.
  \item \textbf{4D representation and evaluation.} The assembled sequence is
    formalized as a 4D texture avatar and evaluated for spatial completeness
    across novel viewpoints.
\end{enumerate}

Stage~1 uses existing methods as components. The contributions of this paper
span Stages~2--4: the cross-topology transfer algorithm, the formal 4D texture
avatar framework, and their integration into an end-to-end pipeline evaluated
at scale.

% ─────────────────────────────────────────────────────────────────────────────
\subsection{Face Reconstruction and Texture Completion}
\label{sec:recon}

\subsubsection{Monocular Geometry Reconstruction}

\medskip\noindent\textbf{Parameter regression.}
Given an input RGB image $\mathbf{I}$, a ResNet-50 encoder regresses a
parameter vector $\Theta = \{\bm{\beta}, \bm{\theta}_p, \bm{\psi}, \bm{\alpha},
\mathbf{l}, \mathbf{c}\}$, where $\bm{\beta} \in \R^{n_{\text{id}}}$ encodes
identity shape, $\bm{\theta}_p \in \R^{6}$ is the 6-DoF head pose,
$\bm{\psi} \in \R^{n_{\text{exp}}}$ encodes per-frame expression, and
$\bm{\alpha}, \mathbf{l}, \mathbf{c}$ capture texture prior, illumination, and
camera parameters respectively.

\medskip\noindent\textbf{Shape decoding.}
Mesh vertices are reconstructed as:
\begin{equation}
  \mathbf{V}(\bm{\beta}, \bm{\psi}) = \bar{\mathbf{S}}
    + \mathbf{B}_{\text{id}}\,\bm{\beta} + \mathbf{B}_{\text{exp}}\,\bm{\psi}
  \;\in\; \R^{N \times 3}
  \label{eq:shape_decode}
\end{equation}
where $\bar{\mathbf{S}}$ is the mean shape and $\mathbf{B}_{\text{id}}$,
$\mathbf{B}_{\text{exp}}$ are the identity and expression basis matrices.
$N$ and $N_f$ denote vertex and face count (values in Section~\ref{sec:exp}).
The mesh carries UV parameterization $\varphi\colon \OUV \to \mathcal{S}$
over $\OUV = [0,1]^2$.

\medskip\noindent\textbf{View-dependent texture baking.}
Each UV texel is colored by projecting the surface point $\varphi(u,v)$ to
image space and reading the pixel value if the point is visible (passes
depth-buffer and back-face tests). Formally:
\begin{equation}
  \Tbaked(u,v) =
  \begin{cases}
    \mathbf{I}(\pi(\varphi(u,v);\bm{\theta}_p))
      & \varphi(u,v) \in \Vis(\bm{\theta}_p) \\
    \text{undefined} & \text{otherwise}
  \end{cases}
  \label{eq:baking}
\end{equation}

\medskip\noindent\textbf{Coverage gap.}
Texture incompleteness is a \emph{deterministic geometric consequence} of the
capture angle, not a stochastic failure. Let $\Oin(\bm{\theta}_p) \subset \OUV$
denote the observable UV region at pose $\bm{\theta}_p$. At a $30^{\circ}$
lateral capture angle:
\begin{equation}
  \tfrac{|\Oin|}{|\OUV|} \approx 0.40\text{--}0.58
  \label{eq:coverage_gap}
\end{equation}
meaning 42--60\% of the UV atlas receives no valid observation regardless of
reconstruction accuracy. For video input, the frame minimizing
$|\theta_{\text{yaw}}|$ is selected to maximize $|\Oin|$; per-frame parameters
$\{\bm{\theta}_{p,t}, \bm{\psi}_t\}_{t=1}^{\tau}$ drive the animation sequence.

% ─────────────────────────────────────────────────────────────────────────────
\subsubsection{UV Texture Completion}

The baked texture $\Tbaked$ and binary validity mask $\cM \in \{0,1\}^{R_u \times R_u}$
(where $R_u$ is the source mesh UV resolution) are passed to a UV-space
inpainting diffusion model~\cite{bai2023uvidm}:
\begin{equation}
  \TIDM = \operatorname{InpaintingModel}(\Tpartial,\; \cM)
  \label{eq:inpainting}
\end{equation}
The model is conditioned on $\Tpartial$ for $\cM=1$ regions and synthesizes
plausible content (skin-tone continuity, approximate bilateral symmetry,
fine surface detail) for $\cM=0$ regions, producing a complete atlas
$\TIDM \in \R^{R_u \times R_u \times 3}$ defined over all of $\OUV$.

\medskip\noindent\textbf{Topology mismatch.}
$\TIDM$ is defined in the source mesh's UV parameterization. The animation mesh
uses a different vertex count, face connectivity, and UV island layout. Because
the same coordinate $(u,v)$ maps to different surface points in each
parameterization, direct UV sampling of $\TIDM$ using animation-mesh UV
coordinates produces incorrect texture. This mismatch requires geometry-guided
transfer described next.

% ─────────────────────────────────────────────────────────────────────────────
\subsection{Cross-Topology Texture Transfer}
\label{sec:transfer}

The completed texture $\TIDM$ lives in the source mesh UV space and must be
rebaked into the target mesh UV space without introducing coverage gaps, seam
artifacts, or topology-induced distortions. We achieve this through a five-step
geometry-aware pipeline.

\subsubsection{Robust Procrustes Alignment}

Let $\mathbf{P}, \mathbf{Q} \in \R^{68 \times 3}$ be the 68 facial landmark
positions on the source and target meshes respectively. Both are centered at
zero mean, and the source is pre-scaled by the ratio of bounding-box diagonals
to eliminate large coordinate-frame offsets between the two meshes.

\medskip\noindent\textbf{Axis-convention detection.}
The two meshes may use different axis conventions (Y-up vs.\ Z-up, etc.).
We test four candidate permutations $\mathcal{A} = \{\texttt{XYZ}, \texttt{XZY},
\texttt{X{-}YZ}, \texttt{XY{-}Z}\}$ and select the best-fit:
\begin{equation}
  A^* = \underset{A \in \mathcal{A}}{\arg\min}\;
        \tfrac{1}{68}\|A(\mathbf{P}_c) - \mathbf{Q}_c\|_F^2
  \label{eq:axis_detect}
\end{equation}
where $\mathbf{P}_c, \mathbf{Q}_c$ are the centered landmark sets.

\medskip\noindent\textbf{Umeyama SVD.}
The rotation $\mathbf{R} \in SO(3)$, scale $s$, and translation $\mathbf{t}$
are recovered from the cross-covariance $\mathbf{K} = \mathbf{Q}_c^\top\mathbf{P}_c / 68$:
\begin{align}
  [\mathbf{U},\bm{\Sigma},\mathbf{V}^\top] &= \operatorname{SVD}(\mathbf{K}),\quad
  \mathbf{D} = \operatorname{diag}(1,1,\det\mathbf{U}\cdot\det\mathbf{V})
  \label{eq:svd} \\
  \mathbf{R} &= \mathbf{U}\mathbf{D}\mathbf{V}^\top,\quad
  s = \tfrac{\operatorname{tr}(\bm{\Sigma}\mathbf{D})}{\operatorname{Var}(\mathbf{P}_c)}
  \label{eq:rotation} \\
  \mathbf{t} &= \bar{\mathbf{q}} - s\,\mathbf{R}\,\bar{\mathbf{p}} \notag
\end{align}
$\mathbf{D}$ is a reflection-correction diagonal preventing the degenerate
$\mathbf{R} = -\mathbf{I}$ solution. Source vertices are then mapped as:
\begin{equation}
  \Vsrc \leftarrow s \cdot \Vsrc \cdot \mathbf{R}^\top + \mathbf{t}
  \label{eq:procrustes_apply}
\end{equation}
yielding a coarse global alignment (landmark RMSE typically 0.003--0.012).

% ─────────────────────────────────────────────────────────────────────────────
\subsubsection{Three-Pass Point-to-Plane ICP Refinement}

Procrustes leaves residual misalignment due to mesh topology differences,
landmark noise, and expression discrepancies. We apply a three-pass ICP
scheme minimizing the point-to-plane energy:
\begin{equation}
  E(\bm{\Xi}) = \sum_{(\mathbf{p},\mathbf{q}) \in \Pi}
    \bigl[(\bm{\Xi}\cdot\mathbf{p} - \mathbf{q}) \cdot \hat{\mathbf{n}}_q\bigr]^2
  \label{eq:icp_obj}
\end{equation}
where $\Pi$ is the inlier correspondence set, $\bm{\Xi} \in SE(3)$ is the
rigid transform, and $\hat{\mathbf{n}}_q$ is the target surface normal.
The three passes are: (1)~coarse P2Plane on voxel-downsampled clouds
($\delta=0.002$, 500 iter, $d=0.05$) to correct large-scale residuals;
(1.5)~brief P2Point (200 iter, $d=0.03$) to stabilize translation;
(2)~fine P2Plane at $\delta/2$ (500 iter, $d=0.02$, convergence $10^{-12}$)
to minimize surface RMSE. An adaptive distance cap
$d_{\text{cap}} = \max(0.03,\, 2\sigma_{\text{surf}})$ is computed from the
post-ICP surface RMSE on 5,000 sampled points.

% ─────────────────────────────────────────────────────────────────────────────
\subsubsection{Vectorized Vertex-Level Color Transfer}
\label{sec:color_transfer}

With the source mesh now co-registered in the target mesh's coordinate space,
we transfer color from $\TIDM$ to each target mesh vertex, producing a
per-vertex color array $\mathbf{C} \in \R^{N \times 3}$ and a confidence weight
$w_i \in [0,1]$ for each vertex $i \in \{1,\dots,N\}$.

For each target vertex $\mathbf{v}_i$, we query its closest point
$(\mathbf{c}_i, d_i, \mathit{tri}_i)$ on the aligned source surface in a
single vectorized call. Rather than hard-gating by distance (which creates
coverage gaps), we assign each vertex a soft confidence weight:
\begin{align}
  w_{\text{dist}}(i) &= \operatorname{clip}(1 - d_i/d_{99},\;0,\;1)
  \label{eq:w_dist} \\
  w_{\text{norm}}(i) &= \operatorname{clip}\!\left(
    \tfrac{\hat{\mathbf{n}}_i \cdot \hat{\mathbf{n}}_{f_i} + 0.8}{1.8},\;0,\;1\right)
  \label{eq:w_norm} \\
  w(i) &= w_{\text{dist}}(i)\cdot w_{\text{norm}}(i)
  \label{eq:w_combined}
\end{align}
where $d_{99}$ is the 99th-percentile distance; $\hat{\mathbf{n}}_i$ and
$\hat{\mathbf{n}}_{f_i}$ are the unit normals at target vertex $i$ and its
closest source triangle $f_i$; the normal term falls to zero at ${\approx}144^\circ$. Since $w(i)>0$ always, every vertex receives a
color (100\% coverage); vertices with low $w(i)$ receive less reliable color
but are never excluded.

The interpolated source UV for vertex $i$ is:
\begin{equation}
  \mathbf{uv}_i = \sum_{k=0}^{2} b_k\,\mathbf{uv}_{\text{src},k}
  \label{eq:bary_weights}
\end{equation}
where $[b_0,b_1,b_2] = \operatorname{barycentric}(\mathbf{c}_i;\,\Vsrc[\mathit{tri}_i])$
are the barycentric coordinates of $\mathbf{c}_i$ within triangle $\mathit{tri}_i$.
and the per-vertex color is sampled bilinearly:
\begin{equation}
  \mathbf{C}_i = \operatorname{BilinearSample}(\TIDM,\; \mathbf{uv}_i)
  \label{eq:bilinear}
\end{equation}
yielding $\mathbf{C} \in \R^{N \times 3}$ with 100\% vertex coverage.

% ─────────────────────────────────────────────────────────────────────────────
\subsubsection{Zero-Loop Vectorized UV Rasterization and Baking}
\label{sec:rast}

The per-vertex color array $\mathbf{C} \in \R^{N \times 3}$ is projected into
the 2D target mesh UV atlas at $R_{\text{tex}} \times R_{\text{tex}}$
resolution ($R_{\text{tex}} = 1024$) via a fully vectorized rasterizer with no
Python-level loops over pixels or faces.

UV coordinates are converted to pixel space; per-triangle AABBs enumerate all
candidate $(f_k, \text{pixel})$ pairs in a flat array without Python loops.
Each candidate is tested with an edge-function barycentric test:
\begin{align}
  \lambda_j &= \tfrac{\text{signed sub-area}_j}{\text{area}(f_k)},
  \quad j\in\{0,1,2\} \notag \\
  \text{inside} &\!=\! (\lambda_0 \!\geq\! 0) \wedge (\lambda_1 \!\geq\! 0)
                        \wedge (\lambda_2 \!\geq\! 0)
  \label{eq:edge_test}
\end{align}
Passing pixels write their face index and barycentric weights $(\lambda_0,\lambda_1,\lambda_2)$
into $\texttt{face\_map}$ and $\texttt{bary\_map}$. Candidates are processed in
chunks of $4\times10^6$ to bound peak RAM.

\medskip\noindent\textbf{Barycentric baking.}
Once \texttt{face\_map} is populated, the final texture color at each covered
texel $(y, x)$ is:
\begin{equation}
  \Tcomplete[y,x] = \sum_{j=0}^{2}\lambda_j\,\mathbf{C}[v_j]
  \label{eq:baking_final}
\end{equation}
where $v_0, v_1, v_2$ are the three target vertex indices of the triangle at
texel $(y, x)$, and $\mathbf{C}$ is the per-vertex color array from
\eqref{eq:bilinear}. This single gather-and-multiply completes in under 0.5~s
for the full $1024\!\times\!1024$ atlas.

\medskip\noindent\textbf{Vectorization speedup.}
Eliminating Python-level loops is critical for practical throughput.
Table~\ref{tab:speedup} reports wall-clock times for three implementation
strategies on the same $N_f = 9{,}976$ face mesh baked to $1024\!\times\!1024$.

\begin{table}[h]
  \caption{UV baking runtime vs.\ implementation strategy ($N_f=9{,}976$,
           $R_{\text{tex}}=1024$, Apple M4 CPU).}
  \label{tab:speedup}
  \centering
  \footnotesize
  \begin{tabular}{@{}lcc@{}}
    \toprule
    Strategy & Time (s) & Speedup \\
    \midrule
    Naïve Python face loop       & ${\approx}300$ & $1\times$ \\
    NumPy loop per face          & ${\approx}45$  & ${\approx}7\times$ \\
    \textbf{Zero-loop flat array (ours)} & ${\approx}8.5$ & ${\approx}35\times$ \\
    \bottomrule
  \end{tabular}
\end{table}

The zero-loop design vectorizes \emph{all} candidate pixel enumeration,
edge-function evaluation, and color gather into flat NumPy arrays, removing
both the face loop and the per-pixel loop. Peak memory is controlled by
chunking candidate pixels at $C_{\max} = 4\,000\,000$, keeping RAM usage
below 4~GB on consumer hardware.

% ─────────────────────────────────────────────────────────────────────────────
\subsubsection{Seam Dilation and Gaussian Boundary Blending}

Let $\mathcal{M}_{\text{cov}}$ be the binary mask of rasterized texels. Three
morphological passes remove remaining artifacts: (A)~a $5\!\times\!5$ dilation
(3 iter) floods single-pixel seam gaps at UV island boundaries; (B)~a
$9\!\times\!9$ dilation (5 iter) fills larger background voids (ears, scalp,
neck); (C)~a Gaussian boundary blend dissolves the hard edge between rasterized
and dilated regions:
\begin{multline}
  \Tout[\mathcal{B}] = \alpha\,\Tout[\mathcal{B}] \\
    + (1-\alpha)\operatorname{GaussBlur}(\Tout,\,\sigma{=}1.5)[\mathcal{B}]
  \label{eq:gauss_blend}
\end{multline}
where $\alpha = \operatorname{GaussBlur}(\mathcal{M}_{\text{cov}}, \sigma{=}3.0)$
and $\mathcal{B} = \{(y,x) : 0.05 < \alpha < 0.95\}$ is the transition band.
The final $\Tcomplete \in \R^{R_{\text{tex}} \times R_{\text{tex}} \times 3}$
has 100\% texel coverage and no visible seam artifacts.

% ─────────────────────────────────────────────────────────────────────────────
\subsection{Formal Definition of the 4D Texture Avatar}

The face surface $\mathcal{S}$ (from §\ref{sec:recon}) carries UV parameterization
$\varphi\colon \OUV \to \mathcal{S}$. A texture map $T\colon \OUV \to \R^3$ assigns
an RGB color to each surface point. Per-frame vertex positions $\mathbf{V}_t$ define
the deformed surface $\mathcal{S}_t$ at frame $t \in [1,\tau]$. A camera with pose
$\theta \in \Theta$ (distinct from capture head pose $\bm{\theta}_p$) projects via
$\pi(\cdot;\theta)\colon \mathcal{S} \to \R^2$, with visible surface set:
\begin{equation}
  \Vis(t,\theta) = \bigl\{p \in \mathcal{S}_t :
    p \text{ visible from } \theta \text{ at frame } t\bigr\}
  \label{eq:vis_def}
\end{equation}

The rendered appearance at pixel $(x,y)$ depends on which texture $T$ is used:
\begin{equation}
  I(x,y,t;\theta) = T\!\left(\varphi^{-1}\!\left(\pi^{-1}(x,y;\theta,t)\right)\right)
  \label{eq:rendered_appearance}
\end{equation}
If $T = \Tpartial$ (output of §\ref{sec:recon}), then $I$ is undefined wherever
$\varphi^{-1}(\pi^{-1}(\cdot)) \notin \Oin$. If $T = \Tcomplete$ (output of
§\ref{sec:transfer}), then $I$ is defined for any $\theta$ and any frame $t$.

\medskip\noindent\textbf{The 4D avatar.}
We define a \emph{4D texture avatar} as the triple:
\begin{equation}
  \cA = \Bigl(\{\mathbf{V}_t\}_{t \in [1,\tau]},\;\; \varphi,\;\; \Tcomplete\Bigr)
  \label{eq:4d_avatar}
\end{equation}
where $\Tcomplete$ is defined on the \emph{entire} UV domain $\OUV$. The descriptor
``4D'' captures that $I(x, y, t; \theta)$ is a function over four free
variables---two spatial screen coordinates $(x, y)$, frame index $t$, and
camera pose $\theta$---yielding a well-defined image for any combination.

\medskip\noindent\textbf{Texture completeness.}
For $I$ to be defined at every $(x,y,t,\theta)$, $T$ must cover all of $\OUV$.
The partial texture from \eqref{eq:baking} satisfies:
\begin{equation}
  \Tpartial(uv) = \begin{cases}
    \mathbf{I}(\pi(\varphi(uv);\bm{\theta}_p)) & uv \in \Oin \\
    \text{undefined} & \text{otherwise}
  \end{cases}
  \label{eq:T_partial_def}
\end{equation}
Since $|\Oin|/|\OUV| \approx 0.40$--$0.58$ (from \eqref{eq:coverage_gap}),
renders at novel poses are systematically undefined. Our completed texture
satisfies the necessary condition:
\begin{equation}
  \Tcomplete(uv) \text{ is defined} \quad \forall\; uv \in \OUV
  \label{eq:T_complete_cond}
\end{equation}
guaranteeing a well-defined $I(x,y,t;\theta)$ at any $(\theta, t)$.

\medskip\noindent\textbf{Concrete representation.}
Each animation frame is the tuple
$f_t = \{\mathbf{V}_t, \mathbf{F}, \mathbf{UV}, \Tcomplete\}$
where $\Tcomplete$ is shared across frames under the identity-constant albedo
assumption. The sequence $\{f_t\}_{t=1}^{\tau}$ is navigable in both viewpoint
and time.

% ─────────────────────────────────────────────────────────────────────────────
\section{Experiments}
\label{sec:exp}

\subsection{Experimental Setup}

We evaluate on \textbf{PolyFace}~\cite{polyface}, a multi-view face dataset
capturing 74 subjects under controlled studio lighting. Each subject is recorded
simultaneously from 7 calibrated cameras at angles C1, C4, C7, C10, C13, C17,
and C24 relative to the frontal axis. We use C4 images (approximately
$30^{\circ}$ lateral offset) as input and C7 images (frontal) as ground truth.
This protocol reflects a realistic capture scenario where the quality of
novel-view rendering is evaluated against the frontal view that most clearly
exposes previously occluded surface regions.

The target animation mesh has $N = 5{,}023$ vertices, $N_f = 9{,}976$ faces,
and $K = 5{,}023$ UV coordinate pairs with a fixed topology shared across all
subjects. The source inpainting mesh has $N_{\text{src}} = 53{,}215$ vertices
with a UV atlas of resolution $R_u = 512$. The output texture atlas is baked at
$R_{\text{tex}} = 1024$.

We report three complementary metrics: \textbf{PSNR} (dB, $\uparrow$) measures
pixel-level reconstruction accuracy; \textbf{SSIM} ($\uparrow$) measures
perceptual structure preservation; and \textbf{LPIPS}~\cite{zhang2018lpips}
($\downarrow$) measures deep feature-level perceptual similarity, correlating
more strongly with human judgment than pixel-level metrics. We compare against
the partial-texture reconstruction baseline~\cite{feng2021deca} without our
texture completion and transfer pipeline.

\subsection{Results}

\noindent\textbf{Quantitative comparison.}
Table~\ref{tab:results} reports per-metric averages and standard deviations
over all 74 subjects.

\begin{table}[h]
  \caption{Quantitative comparison on PolyFace (74 subjects).
           $\uparrow$: higher is better; $\downarrow$: lower is better.}
  \label{tab:results}
  \centering
  \resizebox{\linewidth}{!}{%
  \begin{tabular}{@{}lccc@{}}
    \toprule
    Method & PSNR$\uparrow$ (dB) & SSIM$\uparrow$ & LPIPS$\downarrow$ \\
    \midrule
    Baseline (partial texture) & $13.73 \pm 1.400$ & $0.447 \pm 0.027$ & $0.513 \pm 0.043$ \\
    \textbf{Ours (complete texture)} & $\mathbf{23.80 \pm 1.124}$ & $\mathbf{0.602 \pm 0.054}$ & $\mathbf{0.350 \pm 0.085}$ \\
    $\Delta$ & $\mathbf{+10.07}$ & $\mathbf{+0.155}$ & $\mathbf{-0.163}$ \\
    \bottomrule
  \end{tabular}}
\end{table}

All 74 subjects (100\%) improve on every metric without exception. The
$+10.07$~dB PSNR gain corresponds to reducing mean squared error by a factor
of 10---reflecting the categorical difference between a texture with
structurally blank regions and a fully completed one. The LPIPS improvement of
$-0.163$ confirms that the gain is perceptually meaningful, not merely
numerical. The higher standard deviation of LPIPS (0.085) compared to SSIM
(0.054) suggests that perceptual quality varies more across subjects, likely
due to the inpainting model's sensitivity to hair texture and skin tone
diversity.

\medskip\noindent\textbf{Per-subject consistency.}
The improvement is consistent across subjects: the minimum per-subject PSNR
gain is approximately $+7.5$~dB and the maximum approximately $+14.5$~dB.
Subjects with larger improvements tend to have higher initial occlusion,
confirming that our method provides proportionally greater benefit when more
texture completion is required. This validates that the performance gain is
causally driven by the texture completion mechanism.

\medskip\noindent\textbf{Qualitative comparison.}
Side-by-side renders (Fig.~3) show that the baseline produces sharp geometry
but large blank patches on the ears, temples, and neck. Our method fills these
regions with plausible texture consistent with the visible skin tone, enabling
artifact-free $360^{\circ}$ rendering. The completed texture shows no visible
seam artifacts at the cross-topology transfer boundary, demonstrating the
effectiveness of the ICP-refined alignment and seam dilation post-processing.

\medskip\noindent\textbf{Runtime.}
The cross-topology texture transfer (Procrustes + ICP + rebaking) completes in
approximately 35--55~s per subject on an Apple M4 CPU, dominated by ICP
refinement (${\sim}12$~s) and UV rebaking (${\sim}8$--12~s). The inpainting
model adds 15--30~s. The full pipeline runs in under 90~s on consumer hardware
without GPU acceleration for the transfer stage.

\subsection{Ablation Study}
\label{sec:ablation}

We isolate the contribution of each major design choice in the cross-topology
transfer by removing one component at a time and evaluating on all 74 subjects
from the PolyFace dataset. Table~\ref{tab:ablation} reports mean PSNR, SSIM,
and LPIPS.

\begin{table}[h]
  \caption{Ablation of cross-topology transfer components (74 subjects).}
  \label{tab:ablation}
  \centering
  \resizebox{\linewidth}{!}{%
  \begin{tabular}{@{}lccc@{}}
    \toprule
    Variant & PSNR$\uparrow$ (dB) & SSIM$\uparrow$ & LPIPS$\downarrow$ \\
    \midrule
    w/o soft-confidence (hard gate)     & 18.92 & 0.528 & 0.468 \\
    w/o ICP (Procrustes only)           & 21.37 & 0.562 & 0.412 \\
    w/o Gaussian blending (Pass~C)      & 23.14 & 0.591 & 0.374 \\
    \textbf{Full (ours)}                & \textbf{23.80} & \textbf{0.602} & \textbf{0.350} \\
    \bottomrule
  \end{tabular}}
\end{table}

\medskip\noindent\textbf{Effect of ICP.}
Removing the three-pass ICP and relying on Procrustes alone leaves residual
misalignment—particularly at highly curved regions (ears, nose bridge, jaw
contour) where rigid Procrustes cannot compensate for scale variation
between the two mesh coordinate conventions. The resulting texture shows visible
spatial shift artifacts at these boundaries.

\medskip\noindent\textbf{Effect of soft-confidence weighting.}
Replacing the soft confidence $w(i)$ with a hard distance gate (vertices beyond
a fixed threshold receive zero weight) produces coverage gaps: black or
incomplete texel regions wherever the source and target meshes differ in
surface area coverage. Our soft-confidence formulation avoids all hard
rejections, guaranteeing 100\% texel coverage.

\medskip\noindent\textbf{Effect of Gaussian boundary blending.}
Omitting Pass~C (the Gaussian boundary blend between rasterized and dilated
regions) exposes a visible seam line at UV island boundaries. The hard
transition is perceptually distracting even when geometry alignment is accurate.
The Gaussian blend dissolves this seam within approximately 5 pixels without
altering correctly colored texels.

% ─────────────────────────────────────────────────────────────────────────────
\section{Application Scenarios}

The complete-texture 4D avatar supports free-viewpoint rendering at any
animation time---a property simply absent from any single-viewpoint baked
texture, regardless of resolution or quality. We identify three domains where
this property is a decisive enabling factor.

The most direct application is \emph{personalized avatar creation from a single
photograph}. Because the completed texture covers the entire head surface, the
avatar looks realistic from any observation direction. When the animation
sequence is driven by recorded or synthesized speech, the result is an animated
face avatar whose appearance is angularly consistent---an individualized 4D
representation requiring no multi-view capture hardware and no per-subject
training.

A second domain is \emph{digital cultural and historical preservation}. Museum
collections contain photographs of historical figures where only one image
exists per subject. NeRF-based reconstruction~\cite{guo2021adnerf,park2021nerfface}
requires hundreds of photographs from multiple viewpoints; our pipeline requires
exactly one. The resulting 4D avatar can deliver synthesized speech and be
explored spatially, providing a form of interactive presence that is impossible
from a static photograph.

A third domain is \emph{longitudinal facial geometry and texture monitoring},
relevant to post-surgical recovery tracking and rehabilitation assessment.
A complete 4D reconstruction from monocular video enables low-cost repeated
measurement: the UV texture atlas captures pigmentation changes, scarring
evolution, and skin-tone variation invisible from a fixed frontal camera; the
parametric geometry encodes quantifiable shape changes---facial asymmetry,
volume recovery, jaw mobility---independent of session-to-session head pose
variation.

% ─────────────────────────────────────────────────────────────────────────────
\section{Discussion}

\subsection{The 4D Texture Property: Analysis and Implications}

The 100\% improvement rate is a structural consequence of the problem being
addressed. Partial textures are deterministically incomplete: the blank UV
regions are a direct geometric result of the capture angle, not a stochastic
failure mode. Any completion that fills these regions with plausible content
will necessarily score higher when evaluated against a frontal ground truth that
exposes them.

The deeper significance is quantified by the formal 4D avatar definition in
Section~III-D. The rendered appearance $I(x, y, t; \theta)$ is defined over
the product space $\R^2 \times [1, \tau] \times S^2$, where $S^2$ is the
camera-pose sphere. A partial texture $\Tpartial$ is defined only over the
observed UV subset $\Oin \subset \OUV$---which, for a $30^{\circ}$ lateral
capture angle, covers approximately 40--58\% of the atlas. Concretely, the
fraction of the camera-pose sphere that produces valid renders under
$\Tpartial$ is:
\begin{equation}
  \frac{|\Theta_{\text{valid}}(\Tpartial)|}{|S^2|}
  = \frac{1-\cos\alpha_{\max}}{2} \approx 0.07\text{--}0.12
  \label{eq:valid_sphere}
\end{equation}
for $\alpha_{\max} \approx 30$--$40^{\circ}$. Our complete texture restores
full coverage:
\begin{equation}
  \frac{|\Theta_{\text{valid}}(\Tcomplete)|}{|S^2|} = 1.0
  \label{eq:full_sphere}
\end{equation}

This is a categorical change, not an incremental improvement. A partial-texture
reconstruction is a \emph{2.5D} avatar---animated in time but viewpoint-locked
to a narrow frontal cone covering 7--12\% of the camera sphere. Our
complete-texture avatar is genuinely 4D: navigable across the full time axis
and the full camera-pose sphere simultaneously.

\begin{table}[h]
  \caption{Navigable dimensions of partial-texture vs.\ complete-texture avatars.}
  \label{tab:4d_dims}
  \centering
  \resizebox{\linewidth}{!}{%
  \begin{tabular}{@{}llcc@{}}
    \toprule
    Dimension & Variable & $\Tpartial$ & $\Tcomplete$ \\
    \midrule
    Time           & $t$          & $[1, \tau]$        & $[1, \tau]$  \\
    Camera azimuth & $\theta_{\text{az}}$ & ${\sim}{\pm}30^{\circ}$ & Full $360^{\circ}$ \\
    Camera elevation & $\theta_{\text{el}}$ & ${\sim}{\pm}20^{\circ}$ & Full $180^{\circ}$ \\
    Valid pose coverage & --- & ${\approx}7$--12\% of $S^2$ & 100\% of $S^2$ \\
    \bottomrule
  \end{tabular}}
\end{table}

\medskip\noindent\textbf{Temporal coherence.}
A non-obvious property of our representation is that temporal texture coherence
is guaranteed \emph{by construction} at zero additional cost. Because
$\Tcomplete$ is shared across all frames and the UV parameterization $\varphi$
is fixed regardless of expression, corresponding anatomical points on
consecutive frames always map to identical UV coordinates:
$\varphi^{-1}(p_t) = \varphi^{-1}(p_{t+\Delta t})$, where $p_t$ and
$p_{t+\Delta t}$ are the same anatomical point at consecutive frames. The
texture therefore ``moves with the skin'' automatically. Video-based synthesis
methods must explicitly enforce temporal consistency through recurrent layers or
optical-flow smoothing because they operate in image space where frame-to-frame
correspondence is not encoded.

\medskip\noindent\textbf{Comparison with competing representations.}
Table~\ref{tab:comparison} situates our method relative to the principal
alternative approaches.

\begin{table}[h]
  \caption{Comparison of 4D face representations.}
  \label{tab:comparison}
  \centering
  \resizebox{\linewidth}{!}{%
  \begin{tabular}{@{}lcccc@{}}
    \toprule
    Method & Complete & Free & Per-subject & Single \\
           & texture  & viewpoint & training & photo \\
    \midrule
    2D talking head~\cite{fried2019text} & \texttimes & \texttimes & \texttimes & \checkmark \\
    NeRF-based~\cite{guo2021adnerf,park2021nerfface} & \checkmark & \checkmark & hours & \texttimes \\
    Partial-texture mesh~\cite{feng2021deca} & \texttimes & \texttimes & \texttimes & \checkmark \\
    \textbf{Ours} & \checkmark & \checkmark & \texttimes & \checkmark \\
    \bottomrule
  \end{tabular}}
\end{table}

NeRF-based methods achieve high-quality free-viewpoint renders but require
30--60 minutes of per-subject training on multi-view video. Our system produces
a complete-texture 4D avatar from a single photograph in under 90~s without
per-subject training. The trade-off is that NeRF captures view-dependent
lighting effects---specular reflections, subsurface scattering---that a diffuse
UV texture cannot represent.

\subsection{Limitations and Future Directions}

The primary limitation is \emph{texture hallucination fidelity}. The inpainting
model synthesizes plausible but not necessarily accurate content for occluded
regions. A natural remedy is multi-view supervision: leveraging simultaneous
camera arrays to directly observe held-out UV regions during training, reducing
the fraction of content that must be hallucinated.

A second limitation is the \emph{static texture assumption}. We apply a single
identity texture to all animation frames, which is valid for diffuse skin albedo
but does not capture expression-dependent appearance changes such as wrinkle
deepening or specular highlight shifts. Extending the pipeline to condition the
inpainting model on per-frame expression codes would produce
expression-adaptive texture.

The pipeline also inherits the \emph{geometry accuracy} of the upstream
reconstruction method. Subjects with strong occlusions, extreme head poses, or
atypical face shapes may produce imperfect geometry that the texture transfer
cannot compensate for. Replacing the ICP-based cross-topology alignment with a
learned mesh correspondence network would make the transfer more robust.
Finally, connecting the complete-texture mesh to a speech-driven expression
synthesis module would enable fully generative 4D avatars from text or audio
input alone.

% ─────────────────────────────────────────────────────────────────────────────
\section{Conclusion}

We have presented a complete end-to-end pipeline for free-viewpoint 4D face
avatar reconstruction from a single photograph. The pipeline integrates
monocular face reconstruction, UV-space inpainting, geometry-aware
cross-topology texture transfer, and 4D avatar assembly into a unified system
that produces complete-surface textured animated meshes in under 90~s on
consumer hardware. We formalized the \emph{4D texture avatar} property as a
coverage condition on the UV domain, establishing that partial textures---a
structural limitation of all single-viewpoint monocular methods---violate this
condition for the majority of the camera-pose sphere, while our complete texture
restores it unconditionally. The cross-topology texture transfer module,
combining robust Procrustes alignment, three-pass Point-to-Plane ICP,
soft-confidence vertex color transfer, and zero-loop vectorized UV rebaking,
resolves the topology mismatch between the inpainting mesh and the animation
mesh without coverage gaps or seam artifacts. Quantitative evaluation on 74
subjects demonstrates consistent improvement over the partial-texture baseline
across all subjects and all reported metrics, with the 100\% improvement rate
being a structural consequence of the completeness property. The resulting 4D
avatar is simultaneously navigable across three spatial dimensions and the full
temporal animation axis, enabling downstream applications---personalized avatar
creation, cultural preservation, longitudinal facial monitoring---that require
free-viewpoint rendering from a single photograph.

% ─────────────────────────────────────────────────────────────────────────────
\begin{thebibliography}{16}

\bibitem{feng2021deca}
J.~Feng, H.~Feng, Q.~Li, Z.~Liu, and Z.~Mao,
``DECA: Detailed Expression Capture and Animation,''
\emph{ACM Trans.\ Graph.\ (TOG)}, vol.~40, no.~4, 2021.

\bibitem{li2017flame}
T.~Li, T.~Bolkart, M.~J.~Black, H.~Li, and J.~Romero,
``Learning a Model of Facial Shape and Expression from 4D Scans,''
\emph{ACM Trans.\ Graph.}, vol.~36, no.~6, 2017.

\bibitem{bai2023uvidm}
[UV-IDM authors],
``UV-IDM: Identity-Conditioned Latent Diffusion Model for Face UV-Texture
Generation,''
\emph{arXiv preprint}, 2024.
[Update with correct citation]

\bibitem{blanz1999morphable}
V.~Blanz and T.~Vetter,
``A Morphable Model for the Synthesis of 3D Faces,''
in \emph{Proc.\ SIGGRAPH}, 1999.

\bibitem{paysan2009bfm}
P.~Paysan, R.~Knothe, B.~Amberg, S.~Romdhani, and T.~Vetter,
``A 3D Face Model for Pose and Illumination Invariant Face Recognition,''
in \emph{Proc.\ AVSS}, 2009.

\bibitem{tran2017regressing}
A.~Tran, T.~Hassner, I.~Masi, and G.~Medioni,
``Regressing Robust and Discriminative 3D Morphable Models with a Very Deep
Neural Network,''
in \emph{Proc.\ CVPR}, 2017.

\bibitem{deng2019accurate}
Y.~Deng, J.~Yang, S.~Xu, D.~Chen, Y.~Jia, and X.~Tong,
``Accurate 3D Face Reconstruction with Weakly-Supervised Learning,''
in \emph{Proc.\ CVPRW}, 2019.

\bibitem{danecek2022emoca}
K.~Danecek, M.~J.~Black, and T.~Bolkart,
``EMOCA: Emotion Driven Monocular Face Capture and Animation,''
in \emph{Proc.\ CVPR}, 2022.

\bibitem{deng2018uvgan}
H.~Deng, T.~Cheng, X.~Cao, and Y.~Ma,
``UV-GAN: Adversarial Facial UV Map Completion for Pose-Invariant Face
Recognition,''
in \emph{Proc.\ CVPR}, 2018.

\bibitem{guo2021adnerf}
G.~Guo et~al.,
``AD-NeRF: Audio Driven Neural Radiance Fields for Talking Head Synthesis,''
in \emph{Proc.\ ICCV}, 2021.

\bibitem{park2021nerfface}
K.~Park et~al.,
``NerFace: Deformable Neural Radiance Fields for Dynamic Novel View Synthesis
of Faces,''
in \emph{Proc.\ CVPR}, 2021.

\bibitem{thies2016face2face}
S.~Thies, M.~Zollh\"ofer, M.~Stamminger, C.~Theobalt, and M.~Nie{\ss}ner,
``Face2Face: Real-time Face Capture and Reenactment of RGB Videos,''
in \emph{Proc.\ CVPR}, 2016.

\bibitem{fried2019text}
O.~Fried et~al.,
``Text-based Editing of Talking-head Video,''
\emph{ACM TOG}, vol.~38, no.~4, 2019.

\bibitem{xing2023codetalker}
C.~Xing et~al.,
``CodeTalker: Speech-Driven 3D Facial Animation with Discrete Motion Prior,''
in \emph{Proc.\ CVPR}, 2023.

\bibitem{polyface}
[PolyFace authors],
``PolyFace: A Multi-View Face Dataset for 3D Reconstruction Benchmarking,''
[Conference/Journal, Year].
[Update with correct citation]

\bibitem{zhang2018lpips}
R.~Zhang, P.~Isola, A.~A.~Efros, E.~Shechtman, and O.~Wang,
``The Unreasonable Effectiveness of Deep Features as a Perceptual Metric,''
in \emph{Proc.\ CVPR}, 2018.

\end{thebibliography}

\end{document}
